\documentclass[a4paper,12pt]{article}
\usepackage[utf8]{inputenc}
\usepackage[ngerman]{babel}
\usepackage{amsmath,amssymb}
\usepackage{geometry}\geometry{margin=2.5cm}
\usepackage{enumitem}
\usepackage{booktabs}
\usepackage{tikz}
\usepackage{pgfplots}\pgfplotsset{compat=1.18}
\usepackage{tcolorbox}
\tcbuselibrary{breakable}
\usepackage{xcolor}
\newtcolorbox{begriffe}[1][]{colback=yellow!10, colframe=yellow!60!black, fonttitle=\bfseries, title={Was bedeuten die Begriffe?}, breakable, #1}
\newtcolorbox{wastun}[1][]{colback=orange!5, colframe=orange!60!black, fonttitle=\bfseries, title={Was ist zu tun?}, breakable, #1}
\newcommand{\piv}[1]{{\color{red}\boxed{#1}}}

\title{\"Ubungsblatt 9 -- L\"osungen \\ \large Linear Algebra}
\author{}
\date{}

\begin{document}
\maketitle

%% ============================================================
\subsection*{Aufgabe 35}
%% ============================================================

\textbf{Angabe.} Lineare Abbildung $T \colon \mathbb{R}^5 \to \mathbb{R}^3$,
\[
T(a, b, c, d, e) = (b + c + 4d + 2e,\; a + c + 3d + e,\; 4a - 3b + c - 2e).
\]
\begin{itemize}[nosep]
\item[(a)] Basis und Dimension von $\text{Ker}(T)$.
\item[(b)] Basis und Dimension von $\text{Im}(T)$.
\item[(c)] Allgemeine Form von $v \in \mathbb{R}^5$ mit $T(v) = (1, 2, 5)$.
\item[(d)] Ist $T$ injektiv, surjektiv, bijektiv? (Soll \emph{aus (a),(b) folgen}, ohne neue Rechnung.)
\item[(e)] Ist $T$ invertierbar? Wenn ja: explizite Formel f\"ur $T^{-1}$.
\end{itemize}

\begin{begriffe}
\begin{itemize}[leftmargin=*, nosep]
    \item \textbf{Lineare Abbildung} $T \colon V \to W$: Funktion, die Addition und Skalarmultiplikation respektiert.
    \item \textbf{Kern} $\text{Ker}(T) = \{v \in V : T(v) = 0\}$. Alle Vektoren, die auf den Nullvektor abgebildet werden. Ist ein Unterraum von $V$.
    \item \textbf{Bild} $\text{Im}(T) = \{T(v) : v \in V\} = $ alle erreichbaren Werte. Unterraum von $W$.
    \item \textbf{Darstellungsmatrix} $A$ von $T$: Spalten sind die Bilder der Standard-Einheitsvektoren $T(e_1), T(e_2), \dots$. Dann ist $T(v) = Av$.
    \item \textbf{Injektiv} $\iff \text{Ker}(T) = \{0\} \iff \dim \text{Ker}(T) = 0$.
    \item \textbf{Surjektiv} $\iff \text{Im}(T) = W \iff \dim \text{Im}(T) = \dim W$.
    \item \textbf{Bijektiv} $\iff$ injektiv \emph{und} surjektiv $\iff T$ invertierbar.
    \item \textbf{Dimensionssatz (Rang-Nullit\"ats-Satz)}: $\dim \text{Ker}(T) + \dim \text{Im}(T) = \dim V$ (Definitionsbereich).
\end{itemize}
\end{begriffe}

\begin{wastun}
\begin{enumerate}[leftmargin=*, nosep]
    \item Darstellungsmatrix $A$ aufstellen (Spalten = $T(e_i)$).
    \item Gau\ss-Elimination einmal durchziehen -- damit lassen sich (a), (b), (c) alle direkt ablesen.
    \item \textbf{(a)} Kern: L\"ose homogenes System $Av = 0$. Freie Variablen $\to$ Basisvektoren.
    \item \textbf{(b)} Bild: Pivot-Spalten in der Stufenform markieren $\to$ entsprechende \emph{Original-Spalten} von $A$ sind die Basis.
    \item \textbf{(c)} Partikul\"are L\"osung von $Av = (1,2,5)$ + allgemeine Kernl\"osung.
    \item \textbf{(d)} Aus $\dim \text{Ker}$ und $\dim \text{Im}$ direkt schlie\ss en.
    \item \textbf{(e)} Bijektivit\"at $\iff$ invertierbar.
\end{enumerate}
\end{wastun}

%% --- 35 Darstellungsmatrix ---
\subsubsection*{Schritt 1: Darstellungsmatrix aufstellen}

Bilder der Einheitsvektoren:
\begin{align*}
T(e_1) = T(1,0,0,0,0) &= (0,\; 1,\; 4) \\
T(e_2) = T(0,1,0,0,0) &= (1,\; 0,\; -3) \\
T(e_3) = T(0,0,1,0,0) &= (1,\; 1,\; 1) \\
T(e_4) = T(0,0,0,1,0) &= (4,\; 3,\; 0) \\
T(e_5) = T(0,0,0,0,1) &= (2,\; 1,\; -2)
\end{align*}

Diese f\"unf Vektoren sind die \textbf{Spalten} von $A$:
\[
A = \begin{pmatrix} 0 & 1 & 1 & 4 & 2 \\ 1 & 0 & 1 & 3 & 1 \\ 4 & -3 & 1 & 0 & -2 \end{pmatrix}.
\]

%% --- 35 Stufenform ---
\subsubsection*{Schritt 2: Stufenform von $A$ (einmal f\"ur alle Teilaufgaben)}

\[
\left(\begin{array}{ccccc}
0 & 1 & 1 & 4 & 2 \\
1 & 0 & 1 & 3 & 1 \\
4 & -3 & 1 & 0 & -2
\end{array}\right)
\xrightarrow{Z_1 \leftrightarrow Z_2}
\left(\begin{array}{ccccc}
1 & 0 & 1 & 3 & 1 \\
0 & 1 & 1 & 4 & 2 \\
4 & -3 & 1 & 0 & -2
\end{array}\right)
\xrightarrow{Z_3 \to Z_3 - 4Z_1}
\left(\begin{array}{ccccc}
1 & 0 & 1 & 3 & 1 \\
0 & 1 & 1 & 4 & 2 \\
0 & -3 & -3 & -12 & -6
\end{array}\right)
\]
\[
\xrightarrow{Z_3 \to Z_3 + 3Z_2}
\left(\begin{array}{ccccc}
\piv{1} & 0 & 1 & 3 & 1 \\
0 & \piv{1} & 1 & 4 & 2 \\
0 & 0 & 0 & 0 & 0
\end{array}\right)
\]

Pivots in \textbf{Spalte 1 und 2} (entspricht den Variablen $a, b$). Freie Variablen: $c, d, e$.

%% --- 35 (a) ---
\subsubsection*{(a) Kern von $T$}

L\"ose $Av = 0$. Stufenform $\Rightarrow$ Gleichungen:
\begin{align*}
\text{Zeile 1:} &\quad a + c + 3d + e = 0 \;\Rightarrow\; a = -c - 3d - e \\
\text{Zeile 2:} &\quad b + c + 4d + 2e = 0 \;\Rightarrow\; b = -c - 4d - 2e
\end{align*}

Freie Variablen umbenennen: $c = s$, $d = t$, $e = u$.

\[
(a, b, c, d, e) = (-s - 3t - u,\; -s - 4t - 2u,\; s,\; t,\; u).
\]

Nach $s, t, u$ aufspalten:
\[
\begin{pmatrix} a \\ b \\ c \\ d \\ e \end{pmatrix}
= s \begin{pmatrix} -1 \\ -1 \\ 1 \\ 0 \\ 0 \end{pmatrix}
+ t \begin{pmatrix} -3 \\ -4 \\ 0 \\ 1 \\ 0 \end{pmatrix}
+ u \begin{pmatrix} -1 \\ -2 \\ 0 \\ 0 \\ 1 \end{pmatrix}.
\]

\[
\fbox{\parbox{0.93\textwidth}{\centering
Basis von $\text{Ker}(T)$: $\{(-1,-1,1,0,0),\;(-3,-4,0,1,0),\;(-1,-2,0,0,1)\}$. \\[2pt]
$\dim \text{Ker}(T) = 3.$
}}
\]

%% --- 35 (b) ---
\subsubsection*{(b) Bild von $T$}

$\text{Im}(T) = \text{span}\{T(e_1), T(e_2), T(e_3), T(e_4), T(e_5)\}$, also der Spaltenraum von $A$.

\textbf{Pivot-Spalten der Stufenform markieren $\to$ entsprechende Original-Spalten bilden Basis.} Pivots in Spalten $1, 2 \Rightarrow$ Basis-Spalten sind die \emph{ersten beiden} Spalten der \emph{Original-Matrix} $A$:
\[
\text{Spalte 1 von } A = T(e_1) = (0, 1, 4),\qquad
\text{Spalte 2 von } A = T(e_2) = (1, 0, -3).
\]

\[
\fbox{\parbox{0.93\textwidth}{\centering
Basis von $\text{Im}(T)$: $\{(0, 1, 4),\;(1, 0, -3)\}$. \\[2pt]
$\dim \text{Im}(T) = 2.$
}}
\]

\textbf{Check Dimensionssatz:} $3 + 2 = 5 = \dim \mathbb{R}^5$. \checkmark

%% --- 35 (c) ---
\subsubsection*{(c) Allgemeine $v$ mit $T(v) = (1, 2, 5)$}

\textbf{Idee.} Die Gleichung $T(v) = (1, 2, 5)$ entspricht dem System $Av = b$ mit $b = (1, 2, 5)^\top$. Wir nutzen die \textbf{gleichen Zeilenoperationen} wie in Schritt 2 -- denn die Koeffizientenmatrix $A$ ist dieselbe, nur die rechte Seite ist neu.

\textbf{Schritt 1: Erweiterte Matrix aufstellen und reduzieren.}

\[
\left(\begin{array}{ccccc|c}
0 & 1 & 1 & 4 & 2 & 1 \\
1 & 0 & 1 & 3 & 1 & 2 \\
4 & -3 & 1 & 0 & -2 & 5
\end{array}\right)
\xrightarrow{Z_1 \leftrightarrow Z_2}
\left(\begin{array}{ccccc|c}
1 & 0 & 1 & 3 & 1 & 2 \\
0 & 1 & 1 & 4 & 2 & 1 \\
4 & -3 & 1 & 0 & -2 & 5
\end{array}\right)
\]
\[
\xrightarrow{Z_3 \to Z_3 - 4 Z_1}
\left(\begin{array}{ccccc|c}
1 & 0 & 1 & 3 & 1 & 2 \\
0 & 1 & 1 & 4 & 2 & 1 \\
0 & -3 & -3 & -12 & -6 & -3
\end{array}\right)
\xrightarrow{Z_3 \to Z_3 + 3 Z_2}
\left(\begin{array}{ccccc|c}
\piv{1} & 0 & 1 & 3 & 1 & 2 \\
0 & \piv{1} & 1 & 4 & 2 & 1 \\
0 & 0 & 0 & 0 & 0 & 0
\end{array}\right)
\]

\textbf{Schritt 2: Konsistenz pr\"ufen.} Die letzte Zeile lautet $0 = 0$ -- \textbf{keine} Zeile der Form $(0, 0, 0, 0, 0 \mid \text{etwas} \neq 0)$. Also: System ist konsistent, L\"osungen existieren.

\textbf{Schritt 3: Gleichungen ablesen.} Die zwei Pivot-Zeilen liefern:
\begin{align*}
\text{Zeile 1:} &\quad a + c + 3d + e = 2 \\
\text{Zeile 2:} &\quad b + c + 4d + 2e = 1
\end{align*}
Wie in (a): Pivot-Variablen $a, b$, freie Variablen $c, d, e$.

\textbf{Schritt 4: Partikul\"are L\"osung suchen.} Eine ``partikul\"are L\"osung'' ist \emph{irgendeine} L\"osung -- am einfachsten findet man sie, indem man die freien Variablen auf den simpelsten Wert setzt: $c = d = e = 0$.

\begin{itemize}[leftmargin=*, nosep]
    \item Zeile 1 wird zu: $a + 0 + 0 + 0 = 2 \Rightarrow a = 2$.
    \item Zeile 2 wird zu: $b + 0 + 0 + 0 = 1 \Rightarrow b = 1$.
\end{itemize}
\[
v_0 = (2, 1, 0, 0, 0).
\]

\textbf{Probe:} Setze $v_0$ in $T$ ein:
\[
T(2, 1, 0, 0, 0) = (\underbrace{1 + 0 + 0 + 0}_{b+c+4d+2e},\; \underbrace{2 + 0 + 0 + 0}_{a+c+3d+e},\; \underbrace{8 - 3 + 0 - 0}_{4a-3b+c-2e}) = (1, 2, 5). \;\checkmark
\]

\textbf{Schritt 5: Vom Spezialfall zur allgemeinen L\"osung.}

\textbf{Strukturprinzip:} Wenn $T(v) = b$ eine L\"osung $v_0$ hat, dann ist die \textbf{Menge aller L\"osungen} gegeben durch:
\[
v = v_0 + (\text{beliebiger Vektor aus } \text{Ker}(T)).
\]

\textbf{Warum?} Sei $v$ irgendeine andere L\"osung von $T(v) = b$. Dann ist $T(v - v_0) = T(v) - T(v_0) = b - b = 0$, also liegt $v - v_0$ im Kern, also $v = v_0 + k$ mit $k \in \text{Ker}(T)$.

In (a) haben wir den Kern bestimmt: drei Basisvektoren mit Parametern $s, t, u$. Also:

\[
\fbox{\parbox{0.93\textwidth}{\centering
$v = \underbrace{(2, 1, 0, 0, 0)}_{\text{partikul\"are L\"osung}}
+ \underbrace{s(-1,-1,1,0,0) + t(-3,-4,0,1,0) + u(-1,-2,0,0,1)}_{\text{Kern (siehe (a))}}$, \\[4pt]
$s, t, u \in \mathbb{R}.$
}}
\]

%% --- 35 (d) ---
\subsubsection*{(d) Injektiv / Surjektiv / Bijektiv}

Die Aufgabe sagt: ohne neu zu rechnen, nur aus (a) und (b) ableiten. Daf\"ur brauchen wir drei einfache Regeln:

\begin{itemize}[leftmargin=*, nosep]
    \item \emph{Injektiv} hei\ss t: nur $v = 0$ wird auf $0$ abgebildet. Also: $\dim \text{Ker}(T) = 0$.
    \item \emph{Surjektiv} hei\ss t: jeder Vektor im Zielraum kommt vor. Also: $\dim \text{Im}(T) = \dim W$.
    \item \emph{Bijektiv} hei\ss t: injektiv \emph{und} surjektiv.
\end{itemize}

Jetzt einfach die Zahlen einsetzen:

\textbf{Injektiv?} Aus (a): $\dim \text{Ker}(T) = 3$. Das ist nicht $0$ $\Rightarrow$ \textbf{nicht injektiv}.

\textbf{Surjektiv?} Aus (b): $\dim \text{Im}(T) = 2$. Der Zielraum hat $\dim \mathbb{R}^3 = 3$. Das ist nicht gleich $\Rightarrow$ \textbf{nicht surjektiv}.

\textbf{Bijektiv?} Beides muss gelten -- schon das eine scheitert $\Rightarrow$ \textbf{nicht bijektiv}.

\[
\fbox{\parbox{0.93\textwidth}{\centering
$T$ ist \textbf{weder injektiv, noch surjektiv, noch bijektiv}.
}}
\]

%% --- 35 (e) ---
\subsubsection*{(e) Invertierbar?}

\textbf{Was hei\ss t ``invertierbar''?} $T$ ist invertierbar, wenn es eine Umkehrabbildung $T^{-1}$ gibt, also wenn man $T$ ``r\"uckw\"arts laufen'' lassen kann.

\textbf{Regel:} invertierbar $=$ bijektiv. Beides ist dasselbe.

Aus (d) wissen wir: $T$ ist nicht bijektiv $\Rightarrow$ $T$ ist \textbf{nicht invertierbar}.

\medskip

\textbf{Warum geht das hier prinzipiell nicht?} Eine Umkehrabbildung m\"usste aus $\mathbb{R}^3$ zur\"uck nach $\mathbb{R}^5$ gehen -- aber wir haben nur 3 Zahlen Information und sollen daraus 5 Zahlen rekonstruieren. Das kann nicht eindeutig funktionieren.

\textbf{Kurz gesagt:} $\dim \mathbb{R}^5 = 5 \neq 3 = \dim \mathbb{R}^3$. Schon allein deshalb kann $T$ nicht invertierbar sein -- die Dimensionen passen nicht.

\[
\fbox{\parbox{0.93\textwidth}{\centering
$T$ ist \textbf{nicht invertierbar}.
}}
\]

\bigskip
\hrule
\bigskip

%% ============================================================
\subsection*{Aufgabe 36}
%% ============================================================

\textbf{Angabe.} $T \colon \mathbb{R}^3 \to \mathbb{R}^3$ mit $T(a, b, c) = (3b + c,\; 2a + b,\; a)$.
\begin{itemize}[nosep]
\item[(a)] $T$ invertierbar? Wenn ja: explizite Formel f\"ur $T^{-1}$.
\item[(b)] Aus (a): Kern und Bild bestimmen (einfachste Beschreibungen).
\item[(c)] Aus (a): Injektiv / Surjektiv / Bijektiv?
\end{itemize}

\begin{begriffe}
\begin{itemize}[leftmargin=*, nosep]
    \item \textbf{Invers}: $T^{-1}$ macht $T$ ``r\"uckg\"angig''. F\"ur $(x, y, z) = T(a, b, c)$ findet $T^{-1}$ aus $(x, y, z)$ den Originalvektor $(a, b, c)$.
    \item \textbf{Suchstrategie} f\"ur $T^{-1}$: Gleichungen $T(a, b, c) = (x, y, z)$ nach $a, b, c$ aufl\"osen.
\end{itemize}
\end{begriffe}

%% --- 36 (a) ---
\subsubsection*{(a) Invertierbarkeit und $T^{-1}$}

\textbf{Schritt 1: Gleichungen aufstellen.} $(x, y, z) = T(a, b, c)$ bedeutet:
\begin{align*}
3b + c &= x \\
2a + b &= y \\
a &= z
\end{align*}

\textbf{Schritt 2: Von unten nach oben aufl\"osen.}

\textbf{(a) $a$ aus Gleichung 3:} $a = z$.

\textbf{(b) $b$ aus Gleichung 2:} $b = y - 2a = y - 2z$.

\textbf{(c) $c$ aus Gleichung 1:} $c = x - 3b = x - 3(y - 2z) = x - 3y + 6z$.

\textbf{Eindeutige L\"osung} f\"ur jedes $(x, y, z) \Rightarrow$ $T$ ist invertierbar.

\[
\fbox{\parbox{0.93\textwidth}{\centering
$T^{-1}(x, y, z) = (z,\; y - 2z,\; x - 3y + 6z)$.
}}
\]

\textbf{Probe:} $T(T^{-1}(x,y,z)) = T(z, y-2z, x-3y+6z)$:
\begin{itemize}[nosep, leftmargin=*]
\item 1.\ Komp.: $3(y - 2z) + (x - 3y + 6z) = 3y - 6z + x - 3y + 6z = x$ \checkmark
\item 2.\ Komp.: $2z + (y - 2z) = y$ \checkmark
\item 3.\ Komp.: $z$ \checkmark
\end{itemize}

%% --- 36 (b) ---
\subsubsection*{(b) Kern und Bild}

$T$ ist invertierbar $\Rightarrow$ $T$ ist bijektiv $\Rightarrow$:
\[
\fbox{\parbox{0.93\textwidth}{\centering
$\text{Ker}(T) = \{(0, 0, 0)\}$ \quad (nur der Nullvektor) \\[2pt]
$\text{Im}(T) = \mathbb{R}^3$ \quad (der ganze Zielraum)
}}
\]

%% --- 36 (c) ---
\subsubsection*{(c) Injektiv / Surjektiv / Bijektiv}

$T$ invertierbar $\iff$ $T$ bijektiv $\iff$ $T$ injektiv \emph{und} surjektiv.

\[
\fbox{\parbox{0.93\textwidth}{\centering
$T$ ist \textbf{injektiv, surjektiv, bijektiv}.
}}
\]

\bigskip
\hrule
\bigskip

%% ============================================================
\subsection*{Aufgabe 37}
%% ============================================================

\textbf{Angabe.} $T \colon \mathbb{R}^{2 \times 2} \to \mathbb{R}^{2 \times 2}$ mit $T(X) = M \cdot X$, wobei $M = \begin{pmatrix} 1 & 2 \\ 4 & 8 \end{pmatrix}$.
\begin{itemize}[nosep]
\item[(a)] Basis und Dimension von $\text{Ker}(T)$.
\item[(b)] Basis und Dimension von $\text{Im}(T)$.
\end{itemize}

\begin{begriffe}
\begin{itemize}[leftmargin=*, nosep]
    \item Der Vektorraum ist hier $\mathbb{R}^{2 \times 2}$ (alle $2\times 2$-Matrizen). $\dim \mathbb{R}^{2\times 2} = 4$.
    \item \textbf{Beobachtung:} In $M = \begin{pmatrix} 1 & 2 \\ 4 & 8 \end{pmatrix}$ ist Zeile 2 $= 4 \cdot$ Zeile 1, ebenso Spalte 2 $= 2 \cdot$ Spalte 1. Also hat $M$ Rang $1$ und ist nicht invertierbar.
\end{itemize}
\end{begriffe}

%% --- 37 Berechnung T(X) ---
\subsubsection*{Schritt 1: $T(X)$ ausrechnen mit allgemeinem $X$}

Sei $X = \begin{pmatrix} a & b \\ c & d \end{pmatrix}$. Dann:
\[
T(X) = \begin{pmatrix} 1 & 2 \\ 4 & 8 \end{pmatrix} \begin{pmatrix} a & b \\ c & d \end{pmatrix}
= \begin{pmatrix} a + 2c & b + 2d \\ 4a + 8c & 4b + 8d \end{pmatrix}
= \begin{pmatrix} a + 2c & b + 2d \\ 4(a + 2c) & 4(b + 2d) \end{pmatrix}.
\]

\textbf{Beobachtung:} Zeile 2 von $T(X)$ ist immer das $4$-fache von Zeile 1.

%% --- 37 (a) ---
\subsubsection*{(a) Kern von $T$}

$T(X) = 0 \iff a + 2c = 0$ \emph{und} $b + 2d = 0$ (die Zeile-2-Gleichungen sind automatisch erf\"ullt).

Aus den zwei Gleichungen: $a = -2c$, $b = -2d$. Freie Variablen: $c, d$.

\[
X = \begin{pmatrix} -2c & -2d \\ c & d \end{pmatrix}
= c \underbrace{\begin{pmatrix} -2 & 0 \\ 1 & 0 \end{pmatrix}}_{B_1}
+ d \underbrace{\begin{pmatrix} 0 & -2 \\ 0 & 1 \end{pmatrix}}_{B_2}.
\]

\[
\fbox{\parbox{0.93\textwidth}{\centering
Basis von $\text{Ker}(T)$: $\left\{\begin{pmatrix} -2 & 0 \\ 1 & 0 \end{pmatrix},\; \begin{pmatrix} 0 & -2 \\ 0 & 1 \end{pmatrix}\right\}$. \quad $\dim \text{Ker}(T) = 2$.
}}
\]

%% --- 37 (b) ---
\subsubsection*{(b) Bild von $T$}

Jede Matrix in $\text{Im}(T)$ hat die Form $\begin{pmatrix} \alpha & \beta \\ 4\alpha & 4\beta \end{pmatrix}$ mit $\alpha = a + 2c$, $\beta = b + 2d$. Da $a, c$ (bzw.\ $b, d$) beliebig sind, kann $\alpha$ (bzw.\ $\beta$) jede reelle Zahl annehmen.

\[
\begin{pmatrix} \alpha & \beta \\ 4\alpha & 4\beta \end{pmatrix}
= \alpha \underbrace{\begin{pmatrix} 1 & 0 \\ 4 & 0 \end{pmatrix}}_{C_1}
+ \beta \underbrace{\begin{pmatrix} 0 & 1 \\ 0 & 4 \end{pmatrix}}_{C_2}.
\]

\[
\fbox{\parbox{0.93\textwidth}{\centering
Basis von $\text{Im}(T)$: $\left\{\begin{pmatrix} 1 & 0 \\ 4 & 0 \end{pmatrix},\; \begin{pmatrix} 0 & 1 \\ 0 & 4 \end{pmatrix}\right\}$. \quad $\dim \text{Im}(T) = 2$.
}}
\]

\textbf{Check Dimensionssatz:} $2 + 2 = 4 = \dim \mathbb{R}^{2\times 2}$. \checkmark

\bigskip
\hrule
\bigskip

%% ============================================================
\subsection*{Aufgabe 38}
%% ============================================================

\textbf{Angabe.} Gibt es eine lineare Abbildung mit den gegebenen Eigenschaften? Falls ja: Beispiel. Falls nein: Beweis.
\begin{itemize}[nosep]
\item[(a)] $T \colon \mathbb{R}^2 \to \mathbb{R}^2$ mit $\text{Ker}(T) = \text{Im}(T)$.
\item[(b)] $T \colon \mathbb{R}^3 \to \mathbb{R}^3$ mit $\text{Ker}(T) = \text{Im}(T)$.
\item[(c)] $T \colon \mathbb{R}^6 \to \mathbb{R}^3$ mit $\dim \text{Ker}(T) = 2$.
\end{itemize}

\begin{begriffe}
\begin{itemize}[leftmargin=*, nosep]
    \item \textbf{Schl\"usselformel}: Dimensionssatz $\dim \text{Ker}(T) + \dim \text{Im}(T) = \dim V$ (Definitionsbereich).
    \item \textbf{Trick}: Wenn $\text{Ker}(T) = \text{Im}(T)$ gefordert ist, m\"ussen beide R\"aume insbesondere \textbf{dieselbe Dimension} haben.
    \item \textbf{Trick}: $\dim \text{Im}(T)$ ist immer $\leq \dim W$ (Zielraum).
\end{itemize}
\end{begriffe}

%% --- 38 (a) ---
\subsubsection*{(a) $T \colon \mathbb{R}^2 \to \mathbb{R}^2$ mit $\text{Ker} = \text{Im}$}

\textbf{Dimensionscheck.} Sei $d = \dim \text{Ker} = \dim \text{Im}$ (gleich, weil die R\"aume gleich sein sollen). Dimensionssatz:
\[
d + d = 2 \;\Rightarrow\; d = 1.
\]
Beide R\"aume sind also $1$-dimensional. Das ist konsistent -- es \emph{k\"onnte} ein solches $T$ geben.

\textbf{Beispiel konstruieren.} Idee: $T$ soll alles in einen $1$-dim Unterraum $W$ abbilden ($\to$ $\text{Im} = W$), und $W$ soll genau der Kern sein.

Versuch: $T(a, b) = (b, 0)$.
\begin{itemize}[leftmargin=*, nosep]
    \item $\text{Ker}(T)$: $T(a, b) = (0, 0) \iff b = 0$. Also $\text{Ker}(T) = \{(a, 0) : a \in \mathbb{R}\} = \text{span}\{(1, 0)\}$.
    \item $\text{Im}(T)$: $\{(b, 0) : b \in \mathbb{R}\} = \text{span}\{(1, 0)\}$.
    \item $\text{Ker}(T) = \text{Im}(T)$ \checkmark
\end{itemize}

\[
\fbox{\parbox{0.93\textwidth}{\centering
\textbf{Ja}, z.\,B.\ $T(a, b) = (b, 0)$ erf\"ullt $\text{Ker}(T) = \text{Im}(T) = \text{span}\{(1, 0)\}$.
}}
\]

%% --- 38 (b) ---
\subsubsection*{(b) $T \colon \mathbb{R}^3 \to \mathbb{R}^3$ mit $\text{Ker} = \text{Im}$}

\textbf{Dimensionscheck.} Sei wieder $d = \dim \text{Ker} = \dim \text{Im}$. Dimensionssatz:
\[
d + d = 3 \;\Rightarrow\; 2d = 3 \;\Rightarrow\; d = \tfrac{3}{2}.
\]

$d$ ist die Dimension eines Vektorraums -- muss eine \textbf{ganze Zahl} sein. $\tfrac{3}{2}$ ist es nicht. \textbf{Widerspruch}.

\[
\fbox{\parbox{0.93\textwidth}{\centering
\textbf{Nein}, ein solches $T$ existiert nicht. \\[2pt]
Begr\"undung: aus $\text{Ker} = \text{Im}$ folgt $2 d = 3$, was f\"ur ganzzahliges $d$ unl\"osbar ist.
}}
\]

\textbf{Bemerkung:} Das Argument zeigt allgemein: $\text{Ker}(T) = \text{Im}(T)$ ist nur f\"ur \emph{gerade} Dimensionen des Definitionsbereichs m\"oglich.

%% --- 38 (c) ---
\subsubsection*{(c) $T \colon \mathbb{R}^6 \to \mathbb{R}^3$ mit $\dim \text{Ker}(T) = 2$}

\textbf{Dimensionscheck.} Aus dem Dimensionssatz:
\[
\dim \text{Ker}(T) + \dim \text{Im}(T) = \dim \mathbb{R}^6 = 6 \;\Rightarrow\; \dim \text{Im}(T) = 6 - 2 = 4.
\]

Aber $\text{Im}(T) \subseteq \mathbb{R}^3$, also muss $\dim \text{Im}(T) \leq 3$ sein. Wir h\"atten aber $\dim \text{Im}(T) = 4 > 3$. \textbf{Widerspruch}.

\[
\fbox{\parbox{0.93\textwidth}{\centering
\textbf{Nein}, ein solches $T$ existiert nicht. \\[2pt]
Begr\"undung: aus $\dim \text{Ker} = 2$ folgt $\dim \text{Im} = 4$, aber $\dim \text{Im} \leq \dim \mathbb{R}^3 = 3$.
}}
\]

\end{document}
